\documentclass[11pt]{article}
\usepackage[margin=0.8in]{geometry}
\usepackage{amsmath,amssymb,amsthm}
\usepackage{enumitem}
\usepackage{longtable}
\usepackage{array}
\usepackage{xcolor}
\usepackage{hyperref}
\usepackage{fancyhdr}
\usepackage{titlesec}
\usepackage{tcolorbox}
\usepackage{multicol}
\usepackage{booktabs}
\hypersetup{colorlinks=true, linkcolor=blue!50!black, urlcolor=blue!50!black}
\pagestyle{fancy}
\fancyhf{}
\rhead{\thepage}
\lhead{Engineering Math to Pure Math}
\setlength{\parskip}{0.55em}
\setlength{\parindent}{0pt}
\setlist[itemize]{topsep=0.2em,itemsep=0.2em}
\setlist[enumerate]{topsep=0.2em,itemsep=0.25em}
\newtcolorbox{focusbox}{colback=blue!3!white,colframe=blue!50!black,title=Focus}
\newtcolorbox{intuitionbox}{colback=green!3!white,colframe=green!40!black,title=Intuition}
\newtcolorbox{formalbox}{colback=gray!4!white,colframe=gray!70!black,title=Formal}
\newtcolorbox{walkbox}{colback=orange!5!white,colframe=orange!60!black,title=Walking question}
\newcommand{\N}{\mathbb{N}}
\newcommand{\Z}{\mathbb{Z}}
\newcommand{\Q}{\mathbb{Q}}
\newcommand{\R}{\mathbb{R}}
\newcommand{\C}{\mathbb{C}}
\newcommand{\eps}{\varepsilon}
\newcommand{\pow}{\mathcal{P}}
\theoremstyle{definition}
\newtheorem{definition}{Definition}
\newtheorem{example}{Example}
\newtheorem{exercise}{Exercise}
\theoremstyle{plain}
\newtheorem{theorem}{Theorem}
\newtheorem{lemma}{Lemma}
\theoremstyle{remark}
\newtheorem{remark}{Remark}

\title{\textbf{Start Here: Learning to Think in Proofs}\\A Friendly First Reader Before Velleman}
\author{Engineering Math to Pure Math}
\date{May 2026}
\begin{document}
\maketitle
\tableofcontents
\newpage

\section{Before the Book Opens}

Most people meet mathematics as a long hallway of procedures. Different door, different trick. Derivatives over here. Matrices over there. A binomial expansion in the drawer, a Fourier transform in the attic. Engineering mathematics can be beautiful, but it often rewards speed: recognize the problem, choose the method, compute cleanly, move on.

Proof-based mathematics asks for a different habit. It asks you to slow down enough to know what the words mean.

That sounds small. It is not small.

When a mathematician says ``Let $f:A\to B$ be injective,'' there is a whole little machine sitting inside the sentence. There is a domain $A$, a codomain $B$, a rule $f$, and a promise: two different inputs never collapse to the same output. If you miss any part of that machine, the later proof starts to feel like fog. Not because you are bad at math, but because the sentence was carrying more cargo than it seemed.

\begin{focusbox}
Your first job is not to memorize proof templates. Your first job is to become a careful reader of mathematical sentences.
\end{focusbox}

This reader starts at the floor: statements, variables, quantifiers, sets, functions, and the first proof moves. It is meant to make Velleman's \emph{How to Prove It} feel less like walking into a storm.

\section{Chapter 1: What Is a Mathematical Statement?}

A mathematical statement is a sentence that is true or false. It does not have to be easy to decide. It just has to have a truth value.

\begin{example}
The sentence ``$7$ is prime'' is a statement. It is true.
\end{example}

\begin{example}
The sentence ``$x$ is prime'' is not yet a statement. It depends on $x$. If $x=7$, true. If $x=8$, false. Until $x$ is known, this is a predicate.
\end{example}

\begin{intuitionbox}
In programming terms, a statement is like a Boolean value. A predicate is like a Boolean-valued function waiting for input.
\end{intuitionbox}

This distinction matters because proof writing is full of sentences with moving parts. If you do not know which variables are fixed, arbitrary, or waiting to be chosen, the proof becomes slippery.

\subsection{Statements Versus Commands}

``Compute the derivative of $x^2$'' is not a statement. It is a command. ``The derivative of $x^2$ is $2x$'' is a statement. It is true.

``Find a prime larger than $100$'' is not a statement. ``There exists a prime larger than $100$'' is a statement. It is true.

Mathematics turns questions and commands into statements so they can be proved or disproved.

\subsection{Exercises}

\begin{enumerate}
\item Which are statements?
\begin{enumerate}
\item $3+5=8$.
\item $x+5=8$.
\item There are infinitely many prime numbers.
\item Solve $x^2=9$.
\item Every even integer greater than $2$ is a sum of two primes.
\end{enumerate}
\item Turn each command into a statement:
\begin{enumerate}
\item Find a number whose square is $25$.
\item Show that $n^2+n$ is even.
\item Decide whether $\sqrt2$ is rational.
\end{enumerate}
\end{enumerate}

\section{Chapter 2: The Little Words That Control Everything}

The words ``for every'' and ``there exists'' are small, but they control the entire shape of a proof.

\[
\forall \quad \text{means ``for every.''}
\]
\[
\exists \quad \text{means ``there exists.''}
\]

Compare these two statements:
\[
\forall x\in\R,\ \exists y\in\R \text{ such that } x+y=0.
\]
\[
\exists y\in\R \text{ such that } \forall x\in\R,\ x+y=0.
\]

The first says every real number has an additive opposite. True. If $x=10$, choose $y=-10$. If $x=-\pi$, choose $y=\pi$.

The second says there is one magic number $y$ that cancels every real number. False. No single number cancels both $10$ and $11$.

Same ingredients. Different order. Different universe.

\begin{focusbox}
Quantifier order is not punctuation. It is meaning.
\end{focusbox}

\subsection{How to Read a Universal Statement}

Suppose the claim is:
\[
\forall n\in\Z,\ n^2+n \text{ is even}.
\]

To prove it, you do not get to pick your favorite integer. You must handle an arbitrary integer. The usual proof begins:

\begin{quote}
Let $n\in\Z$ be arbitrary.
\end{quote}

That sentence is doing work. It means: I am not using any secret special property of $n$. Whatever I prove now will apply to all integers.

\subsection{How to Read an Existential Statement}

Suppose the claim is:
\[
\exists n\in\Z \text{ such that } n^2=49.
\]

To prove it, you give a witness. For example, $n=7$.

\begin{quote}
Take $n=7$. Then $n\in\Z$ and $n^2=49$. Therefore such an integer exists.
\end{quote}

Existence proofs can be as simple as pointing to the object. Later they become subtler: sometimes you prove an object exists without knowing its name. That can feel strange at first, but even a simple counting argument can do it.

\subsection{Negating Quantifiers}

Negation is where many beginners first feel the floor tilt.

To negate ``every,'' you need ``some not.''
\[
\neg(\forall x\in A,\ P(x)) \Longleftrightarrow \exists x\in A \text{ such that } \neg P(x).
\]

To negate ``some,'' you need ``none.''
\[
\neg(\exists x\in A,\ P(x)) \Longleftrightarrow \forall x\in A,\ \neg P(x).
\]

If someone says, ``Every server is up,'' you do not disprove it by showing most servers are down. One failed server is enough. If someone says, ``Some server is up,'' you disprove it by showing all servers are down.

\subsection{Exercises}

Negate each statement in clear English and symbols.

\begin{enumerate}
\item Every integer is positive.
\item Some real number has square $2$.
\item Every continuous function is differentiable.
\item There exists a largest prime number.
\item For every user, there exists a password that user remembers.
\end{enumerate}

\section{Chapter 3: If-Then Statements}

Many theorems have the shape:
\[
P\Rightarrow Q.
\]

Read this as: if $P$ is true, then $Q$ is true. The part $P$ is the hypothesis. The part $Q$ is the conclusion.

\begin{example}
If $n$ is divisible by $4$, then $n$ is even.
\end{example}

Here the hypothesis is ``$n$ is divisible by $4$.'' The conclusion is ``$n$ is even.''

A direct proof starts by assuming the hypothesis and ends by proving the conclusion.

\subsection{A First Direct Proof}

\begin{theorem}
If $n\in\Z$ is divisible by $4$, then $n$ is even.
\end{theorem}

\begin{proof}
Suppose $n$ is divisible by $4$. Then there exists $k\in\Z$ such that $n=4k$. Since $4k=2(2k)$ and $2k\in\Z$, there exists an integer $m=2k$ such that $n=2m$. Therefore $n$ is even.
\end{proof}

That proof may feel almost too obvious. Good. The point is not surprise; the point is structure.

Notice what happened:

\begin{enumerate}
\item We assumed the hypothesis.
\item We unpacked the definition of divisible by $4$.
\item We rearranged the expression to match the definition of even.
\item We concluded the target.
\end{enumerate}

\begin{focusbox}
Many beginner proofs are definition unpacking followed by definition repacking.
\end{focusbox}

\section{Chapter 4: Definitions Are Machines}

A definition is not a decorative explanation. It is a machine you can run.

\begin{definition}
An integer $n$ is even if there exists $k\in\Z$ such that $n=2k$.
\end{definition}

\begin{definition}
An integer $n$ is odd if there exists $k\in\Z$ such that $n=2k+1$.
\end{definition}

If a proof says ``$n$ is even,'' you may replace that phrase with ``$n=2k$ for some integer $k$.'' If your goal is to prove ``$n$ is even,'' your job is to produce an integer $k$ such that $n=2k$.

This is the first great proof trick, and it is not really a trick.

\subsection{Worked Example: Odd Times Odd}

\begin{theorem}
If $m$ and $n$ are odd integers, then $mn$ is odd.
\end{theorem}

\begin{proof}
Suppose $m$ and $n$ are odd integers. Then there exist integers $a$ and $b$ such that
\[
m=2a+1,\qquad n=2b+1.
\]
Then
\[
mn=(2a+1)(2b+1)=4ab+2a+2b+1=2(2ab+a+b)+1.
\]
Since $2ab+a+b\in\Z$, this shows $mn$ has the form $2k+1$ for some integer $k$. Therefore $mn$ is odd.
\end{proof}

The proof has a small story inside it. We start with two odd numbers. We open them up using the definition. We multiply. We arrange the result so it looks like the definition of odd again.

\subsection{Exercises}

\begin{enumerate}
\item Prove that the sum of two even integers is even.
\item Prove that the sum of two odd integers is even.
\item Prove that the sum of an even integer and an odd integer is odd.
\item Prove that if $a\mid b$ and $b\mid c$, then $a\mid c$.
\end{enumerate}

\section{Chapter 5: Sets Without Mysticism}

A set is a collection of objects, called elements. The notation $x\in A$ means $x$ is an element of $A$.

Sets are everywhere because they let mathematics talk about collections without caring whether the objects are numbers, points, functions, matrices, symmetries, or open neighborhoods.

\begin{definition}
$A\subseteq B$ means every element of $A$ is also an element of $B$.
\end{definition}

That definition is already a proof plan.

To prove $A\subseteq B$, begin:

\begin{quote}
Let $x\in A$ be arbitrary. We will show $x\in B$.
\end{quote}

\subsection{Set Equality}

To prove $A=B$, prove two inclusions:
\[
A\subseteq B \quad \text{and} \quad B\subseteq A.
\]

This is like showing two APIs accept exactly the same inputs. First no element of $A$ is missing from $B$. Then no element of $B$ is extra.

\subsection{Worked Example}

\begin{theorem}
For all sets $A,B,C$,
\[
A\cap(B\cup C)=(A\cap B)\cup(A\cap C).
\]
\end{theorem}

\begin{proof}
We prove both inclusions.

First, let $x\in A\cap(B\cup C)$. Then $x\in A$ and $x\in B\cup C$. Since $x\in B\cup C$, either $x\in B$ or $x\in C$.

If $x\in B$, then $x\in A\cap B$, so $x\in (A\cap B)\cup(A\cap C)$.

If $x\in C$, then $x\in A\cap C$, so $x\in (A\cap B)\cup(A\cap C)$.

Thus $A\cap(B\cup C)\subseteq (A\cap B)\cup(A\cap C)$.

Now let $x\in (A\cap B)\cup(A\cap C)$. Then either $x\in A\cap B$ or $x\in A\cap C$.

If $x\in A\cap B$, then $x\in A$ and $x\in B$. Since $x\in B$, we have $x\in B\cup C$. Therefore $x\in A\cap(B\cup C)$.

If $x\in A\cap C$, then $x\in A$ and $x\in C$. Since $x\in C$, we have $x\in B\cup C$. Therefore $x\in A\cap(B\cup C)$.

Thus $(A\cap B)\cup(A\cap C)\subseteq A\cap(B\cup C)$. Therefore the two sets are equal.
\end{proof}

\begin{walkbox}
When you see a set identity, ask: what would it mean for one lonely element $x$ to live on the left? What would it mean for the same $x$ to live on the right?
\end{walkbox}

\section{Chapter 6: Functions Are Structured Promises}

A function $f:A\to B$ assigns each element of $A$ exactly one element of $B$.

The phrase ``exactly one'' is doing two jobs:

\begin{itemize}
\item every input gets an output,
\item no input gets two different outputs.
\end{itemize}

The codomain $B$ is part of the function's declared type. The image is the part of $B$ actually hit:
\[
\operatorname{im}(f)=\{f(a):a\in A\}.
\]

\subsection{Injective}

\begin{definition}
$f:A\to B$ is injective if for all $a_1,a_2\in A$,
\[
f(a_1)=f(a_2)\Rightarrow a_1=a_2.
\]
\end{definition}

Injective means the function does not merge two different inputs.

\subsection{Surjective}

\begin{definition}
$f:A\to B$ is surjective if for every $b\in B$, there exists $a\in A$ such that $f(a)=b$.
\end{definition}

Surjective means the function hits every point in the codomain.

\subsection{Worked Example}

\begin{theorem}
If $f:A\to B$ and $g:B\to C$ are injective, then $g\circ f:A\to C$ is injective.
\end{theorem}

\begin{proof}
Let $a_1,a_2\in A$ and suppose $(g\circ f)(a_1)=(g\circ f)(a_2)$. Then
\[
g(f(a_1))=g(f(a_2)).
\]
Since $g$ is injective, $f(a_1)=f(a_2)$. Since $f$ is injective, $a_1=a_2$. Therefore $g\circ f$ is injective.
\end{proof}

The proof walks backward through the composition. Equality at the far end travels through $g$, then through $f$, until it becomes equality of original inputs.

\section{Chapter 7: Relations and the Birth of Quotients}

A relation is a way of saying which objects are connected to which other objects. Equality is a relation. Less-than is a relation. Divisibility is a relation. Having the same remainder after division by $5$ is a relation.

\begin{definition}
A relation $\sim$ on a set $A$ is an equivalence relation if it is:
\begin{itemize}
\item reflexive: $a\sim a$ for all $a\in A$,
\item symmetric: if $a\sim b$, then $b\sim a$,
\item transitive: if $a\sim b$ and $b\sim c$, then $a\sim c$.
\end{itemize}
\end{definition}

Equivalence relations are machines for saying ``these things count as the same for the purpose we care about.''

\subsection{Example: Clock Arithmetic}

On $\Z$, define $a\sim b$ if $a-b$ is divisible by $5$. Then $7\sim 2$ because $7-2=5$. Also $13\sim 3$ because $13-3=10$.

This relation groups integers by remainder:
\[
\ldots,-10,-5,0,5,10,\ldots
\]
\[
\ldots,-9,-4,1,6,11,\ldots
\]
and so on.

This simple idea grows into quotient groups, quotient rings, topology, and geometry. A quotient is what you get when you decide some differences no longer matter.

\begin{walkbox}
In software, when do two different representations count as the same object? Fractions like $1/2$ and $2/4$ are a mathematical version of that question.
\end{walkbox}

\section{Chapter 8: Induction, the Domino Engine}

Induction proves infinitely many statements by proving two things:

\begin{enumerate}
\item the first domino falls,
\item each fallen domino knocks over the next.
\end{enumerate}

\begin{theorem}
For every $n\geq 1$,
\[
1+2+\cdots+n=\frac{n(n+1)}{2}.
\]
\end{theorem}

\begin{proof}
Base case: When $n=1$, the left side is $1$ and the right side is $\frac{1(2)}{2}=1$. So the statement holds for $n=1$.

Induction step: Suppose that for some $k\geq 1$,
\[
1+2+\cdots+k=\frac{k(k+1)}{2}.
\]
We must prove
\[
1+2+\cdots+k+(k+1)=\frac{(k+1)(k+2)}{2}.
\]
Using the induction hypothesis,
\[
1+2+\cdots+k+(k+1)=\frac{k(k+1)}{2}+(k+1).
\]
Factor:
\[
\frac{k(k+1)}{2}+(k+1)=\frac{k(k+1)+2(k+1)}{2}
=\frac{(k+1)(k+2)}{2}.
\]
Thus the statement holds for $k+1$. By induction, it holds for every $n\geq 1$.
\end{proof}

\section{Chapter 9: Counterexamples Are Not Failures}

A counterexample is a single example that disproves a universal claim.

Claim:
\[
\forall a,b\in\Z,\ a\mid b\Rightarrow b\mid a.
\]

Counterexample: $2\mid 4$, but $4\nmid 2$. So the claim is false.

Counterexamples are not cheap shots. They are boundary markers. They tell you where a theorem stops being true.

\subsection{Exercises}

Find counterexamples.

\begin{enumerate}
\item Every function $f:\R\to\R$ is injective.
\item If $x^2=y^2$, then $x=y$.
\item If $A\cup B=A\cup C$, then $B=C$.
\item Every subset of $\R$ has a smallest element.
\end{enumerate}

\section{Chapter 10: Next Subjects}

Once these ideas become ordinary, Velleman becomes much more readable. You will recognize the machinery instead of being surprised by every sentence.

Real analysis will take limits, continuity, and infinite processes and demand exact language. Abstract algebra will take operations and symmetry and ask what survives when the examples change. Topology will ask what continuity really means when distance is no longer the main character. Differential geometry will put calculus on curved spaces. Algebraic geometry will turn polynomial equations into shapes and shapes back into algebra.

That is the larger arc: mathematics keeps finding ways to preserve meaning while changing the setting.

\begin{focusbox}
The beginner-to-hero path is not a leap. It is a stack of small honest moves: read the claim, unpack the definitions, choose the method, justify each step, revise the proof.
\end{focusbox}

\section{A Two-Week Start Plan}

\begin{longtable}{p{0.12\linewidth}p{0.34\linewidth}p{0.44\linewidth}}
\toprule
Day & Work & Output \\
\midrule
1 & Read Chapters 1-2 & Translate 10 statements and negate 5 \\
2 & Re-read quantifiers & Explain $\forall\exists$ versus $\exists\forall$ aloud \\
3 & Read Chapter 3 & Write 3 direct proof skeletons \\
4 & Read Chapter 4 & Prove parity exercises \\
5 & Review definitions & Make flashcards for even, odd, divides, subset \\
6 & Read Chapter 5 & Prove one set identity slowly \\
7 & Rest and walk & Think through the walking questions without paper \\
8 & Read Chapter 6 & Prove composition of injective functions \\
9 & Work function exercises & Find one counterexample involving images \\
10 & Read Chapter 7 & Prove congruence mod $5$ is an equivalence relation \\
11 & Read Chapter 8 & Write the induction proof from memory \\
12 & Read Chapter 9 & Produce four counterexamples \\
13 & Portfolio day & Rewrite your best two proofs cleanly \\
14 & Velleman launch & Open \emph{How to Prove It}, Chapter 1 \\
\bottomrule
\end{longtable}

\end{document}
